\section*{Leitidee}

Funktionen sind das zentrale Objekt der Analysis: Sie beschreiben, wie Gr��en voneinander abh�ngen, und erm�glichen es dadurch, mathematische Zusammenh�nge und reale Prozesse zu modellieren. Dabei kann eine Funktion auf unterschiedliche Weise betrachtet werden -- durch ihre Funktionsgleichung, ihren Graphen oder durch charakteristische Eigenschaften wie Nullstellen, Monotonie und Symmetrie.

Ausgehend vom allgemeinen \textbf{Funktionsbegriff} werden zun�chst grundlegende \textbf{Eigenschaften von Funktionen} untersucht. Anschlie�end stehen wichtige Funktionsklassen im Mittelpunkt: lineare Funktionen, Potenz- und Wurzelfunktionen, Polynome, rationale Funktionen, Exponential- und Logarithmusfunktionen sowie trigonometrische Funktionen. Jede dieser Funktionsklassen besitzt typische Eigenschaften und eignet sich zur Beschreibung bestimmter mathematischer Zusammenh�nge.

Ein weiterer Schwerpunkt liegt auf dem L�sen von \textbf{Gleichungen}, die aus diesen Funktionen entstehen. Dabei wird deutlich, wie die algebraischen Methoden aus dem vorherigen Kapitel mit der Untersuchung von Funktionen zusammenwirken.

Schlie�lich k�nnen Funktionen miteinander \textbf{verkn�pft und komponiert} werden. Auf diese Weise entstehen aus elementaren Funktionen komplexere Modelle, mit denen sich beispielsweise Wachstumsprozesse oder Schwingungen mathematisch beschreiben lassen.



